\documentclass[11pt,a4paper]{article}
\usepackage[top=1.7cm,bottom=1.5cm,left=1.5cm,right=1.5cm,headsep=8pt,footskip=16pt]{geometry}
\usepackage{enumitem}
\usepackage{amssymb}
\usepackage{array}
\usepackage[table]{xcolor}
\usepackage{fancyhdr}
\usepackage{lastpage}
\usepackage{multicol}
\usepackage{polyglossia}
\setdefaultlanguage[calendar=gregorian, numerals=maghrib]{arabic}
\setotherlanguage{english}
\IfFontExistsTF{Amiri}{%
  \newfontfamily\arabicfont[Script=Arabic, Scale=1.08]{Amiri}
}{%
  \newfontfamily\arabicfont[Script=Arabic, Scale=1.08, Path=../fonts/, Extension=.ttf,
    UprightFont=*-Regular, BoldFont=*-Bold, ItalicFont=*-Italic,
    BoldItalicFont=*-BoldItalic]{Amiri}
}
\IfFontExistsTF{AlBattar}{%
  \newfontfamily\khatt[Script=Arabic, Scale=1.15]{AlBattar}
}{%
  \newfontfamily\khatt[Script=Arabic, Scale=1.15, Path=../fonts/,
    BoldFont=ae_AlBattar.ttf]{ae_AlBattar.ttf}
}
\renewcommand{\headrulewidth}{0.8pt}
\renewcommand{\footrulewidth}{0.8pt}
\fancyhf{}
\rhead{\khatt المدرسة العليا للأساتذة ببوزريعة}
\lhead{\khatt ملحقة بومرداس}
\lfoot{\khatt الأستاذ: د. بن عبد الرزاق يوسف}
\cfoot{\thepage/\pageref{LastPage}}
\pagestyle{fancy}
\parindent 0pt
\setlength{\columnsep}{0.6cm}
\setlength{\columnseprule}{0.4pt}
\setlist[enumerate,1]{label=\textbf{\arabic*.}, leftmargin=*, itemsep=5pt, topsep=2pt}
\begin{document}

\linespread{1.08}\selectfont
\begin{center}
{\LARGE\khatt\textbf{الاختبار الثامن والعشرون في الذكاء الاصطناعي}}\\[4pt]
{\normalsize وضعيات إشكالية ــ أسئلة متعددة الاختيارات ــ الفصول من 1 إلى 6}
\end{center}
\vspace{-2pt}\rule{\textwidth}{0.6pt}\\[2pt]
\textbf{التعليمة:} يتضمن الاختبار خمس وضعيات إشكالية. اقرأ كل وضعية بعناية، ثم اختر التحليل أو التصرف الأنسب
من بين الاختيارات الخمسة؛ لكل وضعية إجابة صحيحة واحدة فقط.
اكتب \textbf{حرف} الإجابة الصحيحة في جدول الإجابات في آخر الاختبار.
لكل وضعية 4 نقاط (المجموع 20 نقطة).
\vspace{4pt}

\noindent
{\large\khatt\textbf{الوضعية 1}}\hfill{\small(الفصل 1)}\par\nopagebreak\vspace{3pt}
\setlength{\fboxsep}{7pt}\colorbox{gray!10}{\parbox{\dimexpr\linewidth-14pt}{يبحث طالب ماستر عن إشكالية لمذكرته في تعليمية العربية، فطلب من أداة توليدية «اقترح إشكالية». جاءت عشر إشكاليات عامة، بعضها مدروس كثيرًا، وبعضها يتعذر إنجازه في سنة واحدة. ويريد استعمال الأداة بطريقة تساعده دون أن تفكر عنه. وقد سمع أن الأداة تقترح إشكاليات «جاهزة للإنجاز». ما الاستعمال الأنسب؟}}\par\nopagebreak
\begin{itemize}[leftmargin=1.0cm, labelsep=6pt, itemsep=3pt, topsep=4pt, parsep=0pt]
  \item[\textbf{أ)}] يعتمد أول إشكالية من القائمة، لأن الأداة ترتّب اقتراحاتها حسب الجودة دائمًا، والإشكالية الأولى هي الأجدر بالبحث، فلا حاجة إلى مراجعة الدراسات السابقة أو استشارة المشرف.
  \item[\textbf{ب)}] يرفض الأداة كليًا في مرحلة اختيار الإشكالية، لأن أي استعانة بها في هذه المرحلة تجعل المذكرة كلها من إنتاجها، ولو صاغ الإشكالية بنفسه وناقشها مع مشرفه.
  \item[\textbf{ج)}] يطلب من الأداة إشكالية «جديدة تمامًا لم يدرسها أحد»، ويعتمدها دون تحقق، لأن الأداة تعرف كل الدراسات السابقة، فإذا قالت إن الإشكالية جديدة فهي جديدة بالضرورة.
  \item[\textbf{د)}] يحدد اهتمامه وميدانه والإمكانات المتاحة، ويستعمل الأداة شريكًا في العصف الذهني ونقد صياغاته، ثم يتحقق من جِدّة الإشكالية في الدراسات السابقة، ويناقشها مع مشرفه قبل اعتمادها.
  \item[\textbf{هـ)}] يطلب مئة إشكالية ويختار أطولها صياغة، لأن طول الإشكالية دليل على عمقها، ثم يعتمدها مباشرة دون أن يتحقق من جِدّتها أو من إمكانية إنجازها في سنة واحدة.
\end{itemize}
\par\vspace{6pt}
\noindent
{\large\khatt\textbf{الوضعية 2}}\hfill{\small(الفصل 2)}\par\nopagebreak\vspace{3pt}
\setlength{\fboxsep}{7pt}\colorbox{gray!10}{\parbox{\dimexpr\linewidth-14pt}{درّب فريق نموذجًا يتعرف على أخطاء الكتابة في إنشاءات جُمعت كلها من متوسطات في وسط مدينة كبيرة، ثم قاسوا دقته على عينة أخرى من المتوسطات نفسها فكانت مرتفعة. ويريدون تعميمه على متوسطات الولاية كلها، بما فيها المناطق الريفية. ومديرية التربية مستعدة لتمويل التعميم إذا قُدّم لها تقرير مقنع. ما التقدير الأدق؟}}\par\nopagebreak
\begin{itemize}[leftmargin=1.0cm, labelsep=6pt, itemsep=3pt, topsep=4pt, parsep=0pt]
  \item[\textbf{أ)}] يُرفض النموذج نهائيًا، لأن أي نموذج تدرّب على منطقة واحدة لا يمكن تحسينه أو اختباره على مناطق أخرى، والحل الوحيد بناء نموذج جديد من الصفر لكل متوسطة على حدة.
  \item[\textbf{ب)}] الدقة المقيسة لا تضمن الأداء في المناطق الريفية، لأن عينة الاختبار من المجتمع نفسه الذي جاءت منه بيانات التدريب؛ فيُختبر النموذج على عينة تمثّل المناطق المستهدفة، ويُكمَّل تدريبه عند الحاجة.
  \item[\textbf{ج)}] يُعمَّم النموذج ويُقاس أداؤه بعد سنة من الاستعمال، لأن الاختبار المسبق على المناطق الريفية مضيعة للوقت، والأخطاء ستظهر وتُصحَّح تلقائيًا أثناء الاستعمال الفعلي في الأقسام.
  \item[\textbf{د)}] الدقة المرتفعة تكفي لتعميم النموذج، لأن الأخطاء الكتابية واحدة عند كل التلاميذ في كل المناطق، فلا حاجة إلى أي اختبار إضافي على عينة من المناطق الريفية قبل التعميم، ويوفر على الفريق تكاليف التجريب.
  \item[\textbf{هـ)}] يُختبر النموذج مرة ثانية على عينة جديدة من متوسطات المدينة نفسها، فإن بقيت الدقة مرتفعة عُمِّم على الولاية كلها، لأن ثبات الدقة دليل على صلاحيته في كل المناطق، فيطمئن الفريق قبل التعميم.
\end{itemize}
\par\vspace{6pt}
\noindent
{\large\khatt\textbf{الوضعية 3}}\hfill{\small(الفصل 3)}\par\nopagebreak\vspace{3pt}
\setlength{\fboxsep}{7pt}\colorbox{gray!10}{\parbox{\dimexpr\linewidth-14pt}{طلبت أستاذة من الأداة «أمثلة على الأخطاء الشائعة في كتابة الهمزة» لنشاط تصحيحي، فجاءت أمثلة صحيحة الكتابة أحيانًا، وأخرى خاطئة دون بيان وجه الخطأ. وتريد أمثلة خاطئة فعلًا، يقابل كلًّا منها التصويب والتعليل، بمستوى السنة الثالثة متوسط. وستستعمل الأمثلة في ثلاثة أقسام، وتعيد استعمال الطلب في دروس أخرى. أي طلب أنسب؟}}\par\nopagebreak
\begin{itemize}[leftmargin=1.0cm, labelsep=6pt, itemsep=3pt, topsep=4pt, parsep=0pt]
  \item[\textbf{أ)}] تقدّم مثالين نموذجيين وتطلب عشرة على منوالهما، ثم تعتمد الأمثلة دون مراجعة، لأن النماذج تضمن وحدها صحة التصويب والتعليل في كل الأمثلة المولَّدة دون استثناء.
  \item[\textbf{ب)}] تطلب أمثلة خاطئة من «كتابات التلاميذ الحقيقية في الجزائر»، لأن الأداة تملك نسخًا من إنشاءات التلاميذ في كل المدارس، فتأتي الأمثلة أصيلة دون حاجة إلى نماذج أو مراجعة، فتكون أقرب إلى واقع التلاميذ.
  \item[\textbf{ج)}] تكتب «أمثلة كثيرة على أخطاء الهمزة» فقط، لأن الأداة تميّز تلقائيًا بين الكتابة الخاطئة والصحيحة، وتضيف التصويب والتعليل دون طلب، فلا حاجة إلى نموذج أو مراجعة، ويسهل استعمالها مع كل الأقسام.
  \item[\textbf{د)}] تطلب أمثلة صحيحة الكتابة فقط، ثم تطلب من التلاميذ اكتشاف الأخطاء فيها، لأن الأمثلة الصحيحة تكفي في النشاط التصحيحي، ولا حاجة إلى أمثلة خاطئة أو تعليلات مكتوبة.
  \item[\textbf{هـ)}] تقدّم مثالين نموذجيين بالشكل المطلوب: الجملة الخاطئة، ثم التصويب، ثم القاعدة في سطر، وتطلب عشرة أمثلة على منوالهما في الهمزة المتوسطة والمتطرفة، ثم تراجع كل مثال قبل استعماله.
\end{itemize}
\par\vspace{6pt}
\noindent
{\large\khatt\textbf{الوضعية 4}}\hfill{\small(الفصل 4)}\par\nopagebreak\vspace{3pt}
\setlength{\fboxsep}{7pt}\colorbox{gray!10}{\parbox{\dimexpr\linewidth-14pt}{نشر أستاذ مقالًا في مجلة علمية محكّمة كتب أغلب فقراته أداة ذكاء اصطناعي، دون إفصاح، مع أن سياسة المجلة تشترط التصريح بكل استعمال. واكتشف أحد المراجعين ذلك بعد النشر، وطلبت المجلة توضيحًا. ويرى الأستاذ أن الأفكار أفكاره. ويطالب بعض زملائه بأن يدافع عن نفسه بالإنكار. ما التقدير الأدق؟}}\par\nopagebreak
\begin{itemize}[leftmargin=1.0cm, labelsep=6pt, itemsep=3pt, topsep=4pt, parsep=0pt]
  \item[\textbf{أ)}] مخالفة صريحة لسياسة المجلة ولمبدأ الشفافية، فأصالة الأفكار لا تُسقط واجب الإفصاح عن دور الأداة في الكتابة؛ فيقدّم توضيحًا صادقًا، ويتحمل ما تقرره المجلة من تصحيح أو سحب.
  \item[\textbf{ب)}] يُنكر الأستاذ استعمال الأداة، لأن برامج الكشف غير دقيقة، ولا يمكن للمجلة إثبات الاستعمال، والإنكار حق مشروع ما دامت الأفكار الأساسية من إنتاجه الشخصي.
  \item[\textbf{ج)}] المخالفة بسيطة، فيكفي أن يرسل للمجلة قائمة بالأدوات التي استعملها دون أي توضيح آخر، لأن ذكر اسم الأداة بعد النشر يعادل الإفصاح المسبق تمامًا في كل المجلات، ويُغلق الملف بذلك.
  \item[\textbf{د)}] لا مخالفة، لأن الأفكار أفكاره، والصياغة مسألة شكلية لا تدخل في الأمانة العلمية، وسياسات المجلات حول الذكاء الاصطناعي توصيات غير ملزمة لا تترتب عليها أي إجراءات.
  \item[\textbf{هـ)}] المسؤولية على المراجعين الذين لم يكتشفوا الاستعمال قبل النشر، لا على الأستاذ، لأن مهمة الكشف عن الذكاء الاصطناعي تقع على المجلة وحدها، فلا يلزمه أي توضيح أو تصحيح، فينتهي الأمر دون أي إجراء.
\end{itemize}
\par\vspace{6pt}
\noindent
{\large\khatt\textbf{الوضعية 5}}\hfill{\small(الفصل 6)}\par\nopagebreak\vspace{3pt}
\setlength{\fboxsep}{7pt}\colorbox{gray!10}{\parbox{\dimexpr\linewidth-14pt}{طلبت أستاذة من أداة توليد صور رسم «شجرة إعرابية» لجملة فعلية طويلة، لتوضيح العلاقات بين عناصرها. جاء الرسم أنيقًا، لكن بعض الأسهم تربط الكلمات بعوامل خاطئة، وبعض الكلمات مكتوبة بحروف مشوهة، ورتبة المفعول به مقلوبة. وتريد عرض المخطط على السبورة الذكية وطبعه في المطويات. ما الخطة الأنسب لإنتاج مخطط سليم؟}}\par\nopagebreak
\begin{itemize}[leftmargin=1.0cm, labelsep=6pt, itemsep=3pt, topsep=4pt, parsep=0pt]
  \item[\textbf{أ)}] تعتمد الرسم كما هو، لأن الأناقة البصرية تجذب التلاميذ، والأخطاء في الأسهم والرتبة طفيفة لا تؤثر في فهم العلاقات الإعرابية، والتلاميذ يصححون الحروف المشوهة بأنفسهم.
  \item[\textbf{ب)}] تطلب من الأداة نفسها شرح الرسم، فإن جاء الشرح صحيحًا اعتمدت الرسم، لأن صحة الشرح تثبت صحة الأسهم والرتبة، حتى لو بدت مخالفة للعلاقات الإعرابية في الجملة، فتكسب وقتًا في إعداد الدرس.
  \item[\textbf{ج)}] تحدد العلاقات الإعرابية بنفسها، ثم تستعمل أداة مخططات أو برنامج رسم تضبط فيه العقد والأسهم يدويًا، أو تطلب من أداة نصية شيفرة مخطط تراجعها، لأن مولّد الصور لا يضمن دقة البنية.
  \item[\textbf{د)}] تكرر الطلب حتى يظهر رسم بلا أخطاء ظاهرة، ثم تعتمده دون التحقق من كل سهم، لأن الرسم الخالي من الأخطاء الظاهرة يكون سليم البنية الإعرابية بالضرورة، فيأتي المخطط جاهزًا للعرض.
  \item[\textbf{هـ)}] تتخلى عن أي مخطط إعرابي، لأن المخططات تربك التلاميذ دائمًا، ولا يمكن رسمها بدقة بأي أداة رقمية، حتى ببرامج الرسم التي تُضبط فيها العقد والأسهم يدويًا.
\end{itemize}
\par\vspace{6pt}
\begin{center}
{\khatt\textbf{جدول الإجابات}}\\[3pt]
\renewcommand{\arraystretch}{1.6}
\begin{tabular}{|>{\centering\arraybackslash}m{1.9cm}|*{5}{>{\centering\arraybackslash}m{1.2cm}|}}
\hline
\textbf{رقم السؤال} & \textbf{1} & \textbf{2} & \textbf{3} & \textbf{4} & \textbf{5} \\ \hline
\textbf{الإجابة} &  &  &  &  &  \\[8pt] \hline
\end{tabular}

\vspace{8pt}
\renewcommand{\arraystretch}{2.2}
\begin{tabular}{|>{\centering\arraybackslash}m{7.5cm}|>{\centering\arraybackslash}m{3.2cm}|>{\centering\arraybackslash}m{3.2cm}|}
\hline
\textbf{الاسم واللقب} & \textbf{الفوج} & \textbf{العلامة} \\ \hline
 & & \hfill /20 \\ \hline
\end{tabular}
\end{center}
\end{document}
